\documentclass[12pt,a4paper]{article}

\usepackage{fontawesome5}
\providecommand{\faPencilAlt}{\faIcon{pencil-alt}}

\usepackage{styles/couleurs}
\usepackage{styles/mise_en_page}
\usepackage{styles/boites}
\usepackage{styles/maths}
\usepackage{styles/figures}

\usepackage{enumitem}
\usepackage{array}
\usepackage{tabularx}
\usepackage{amssymb}
\usepackage{pdflscape}
\usepackage{tikz}
\usepackage{xparse}

\setlist[itemize]{leftmargin=1.4em,itemsep=0.15em,topsep=0.2em}
\setlist[enumerate]{leftmargin=*,itemsep=0.2em,topsep=0.2em}

\newcommand{\lignepointillee}{\noindent\dotfill\par\vspace{0.48em}}

\newcommand{\zonelignes}[1]{%
  \begin{tcolorbox}[
    enhanced,
    colback=white,
    colframe=NeutreC!40,
    boxrule=0.5pt,
    arc=1.5mm,
    left=3mm,
    right=3mm,
    top=2mm,
    bottom=1mm,
    before skip=0pt,
    after skip=0pt
  ]
    \foreach \i in {1,...,#1}{\lignepointillee}
  \end{tcolorbox}
}

\newcounter{exohc}

\NewDocumentCommand{\exohc}{m +m}{%
  \refstepcounter{exohc}
  \par\medskip
  \noindent
  \begin{minipage}[t]{0.47\linewidth}
    \vspace{0pt}
    {\sffamily\bfseries\color{TitreC}$\square$ Exercice \theexohc}
    \vspace{0.2em}
    #2
  \end{minipage}
  \hfill
  \begin{minipage}[t]{0.50\linewidth}
    \vspace{0pt}
    \zonelignes{#1}
  \end{minipage}
  \par\medskip
}

\begin{document}
\begin{landscape}

\begin{center}
  {\sffamily\LARGE\bfseries\color{TitreC}Exercices hors classe}\\[0.2em]
  {\sffamily\large\bfseries\color{ActionC}Statistiques et proportionnalité}
\end{center}

\remarque{Cette fiche est donnée à l'avance.
Ne faire que les exercices indiqués par le professeur.}

\section*{A. Reconnaître une situation de proportionnalité}

\exohc{5}{
Indique si la situation est proportionnelle. Justifie.

\begin{enumerate}[label=\alph*)]
  \item Nombre de places de cinéma identiques et prix total.
  \item Âge d'un enfant et âge de sa mère.
  \item Longueur du côté d'un carré et périmètre.
  \item Longueur du côté d'un carré et aire.
\end{enumerate}
}

\exohc{5}{
Une recette utilise \(600\) g de farine pour \(8\) personnes.

\begin{enumerate}[label=\alph*)]
  \item Quelles sont les deux grandeurs ?
  \item Quelles sont leurs unités ?
  \item La situation est-elle proportionnelle ?
\end{enumerate}
}

\section*{B. Compléter des tableaux}

\exohc{5}{
Complète.

\[
\begin{array}{c|cccc}
\text{Nombre d'objets} & 3 & 5 & 8 & 12\\
\hline
\text{Prix en €} & 7{,}5 & & 20 &
\end{array}
\]
}

\exohc{6}{
Pour \(4\) kg de pommes, on paie \(10\) €.

\[
\begin{array}{c|ccccc}
\text{Masse en kg} & 1 & 2 & 3{,}5 & 4 & 7\\
\hline
\text{Prix en €} & & & & 10 &
\end{array}
\]

Complète en utilisant le retour à l'unité.
}

\exohc{5}{
Complète en utilisant l'additivité.

\[
\begin{array}{c|cccc}
\text{Durée en h} & 2 & 3 & 5 & 8\\
\hline
\text{Distance en km} & 90 & 135 & & 
\end{array}
\]
}

\section*{C. Résoudre de petits problèmes}

\exohc{6}{
Un rôti coûte \(20\) € par kilogramme.

\begin{enumerate}[label=\alph*)]
  \item Quelle masse coûte \(30\) € ?
  \item Quelle masse coûte \(18\) € ?
  \item Quelle masse coûte \(37\) € ?
\end{enumerate}
}

\exohc{6}{
Sur un plan, \(4\) cm correspondent à \(290\) m dans la réalité.

\begin{enumerate}[label=\alph*)]
  \item À quelle distance réelle correspondent \(20\) cm ?
  \item À quelle longueur sur le plan correspondent \(725\) m ?
\end{enumerate}
}

\section*{D. Effectifs, fréquences et pourcentages}

\exohc{6}{
Dans un club de \(50\) membres, \(20\) pratiquent le tennis,
\(15\) le badminton, \(10\) le tennis de table et \(5\) le squash.

Complète un tableau avec les effectifs, les fréquences et les pourcentages.
}

\exohc{5}{
Dans une classe de \(25\) élèves, \(10\) viennent à pied.

\begin{enumerate}[label=\alph*)]
  \item Quelle est la fréquence ?
  \item Quel est le pourcentage ?
  \item Combien d'élèves ne viennent pas à pied ?
\end{enumerate}
}

\section*{E. Diagrammes circulaires}

\exohc{6}{
Complète.

\[
\begin{array}{c|cccc}
\text{Pourcentage} & 50\% & 25\% & 20\% & 5\%\\
\hline
\text{Angle} & & & &
\end{array}
\]
}

\exohc{7}{
Une enquête auprès de \(40\) élèves donne :

\[
\begin{array}{c|cccc}
\text{Choix} & A & B & C & D\\
\hline
\text{Effectif} & 12 & 10 & 8 & 10
\end{array}
\]

Calcule le pourcentage et l'angle associés à chaque choix.
}

\exohc{5}{
Un diagramme annonce \(35\%\), \(28\%\), \(22\%\) et \(20\%\).

\begin{enumerate}[label=\alph*)]
  \item Peut-il représenter une répartition complète ?
  \item Quelle vérification as-tu effectuée ?
\end{enumerate}
}

\end{landscape}
\end{document}
